\documentclass[10pt,a4paper]{amsart}
\usepackage[margin=19mm]{geometry}
\usepackage{amsmath,amssymb,mathtools}
\usepackage[scr=rsfs,cal=euler]{mathalfa}
\usepackage{xfrac}
\usepackage[unicode,hidelinks,bookmarks=false]{hyperref}
\newcommand{\ie}{i.e.\ }
\newcommand{\cf}{cf.\ }
\newcommand{\resp}{respectively\ }
\newcommand{\Subsection}{\subsection}
\providecommand{\nobreakdash}{\nobreak\hbox{-}}
\setlength{\emergencystretch}{2em}
\multlinegap0pt
\usepackage{bm}
\usepackage{bbm}
\usepackage{dsfont}
\usepackage{xspace}
\usepackage{cancel}
\usepackage{mathtools}
\usepackage{amsthm}
\makeatletter
\@ifundefined{proposition}{\newtheorem{proposition}{Proposition}}{}
\makeatother
\makeatletter
\newcommand {\CT} {\operatorname {CT}}
\expandafter\ifx\csname conjecture\endcsname\relax
\newenvironment{conjecture}{\begin{conj}}{\end{conj}}
\fi
\expandafter\ifx\csname corollary\endcsname\relax
\newenvironment{corollary}{\begin{coro}}{\end{coro}}
\fi
\expandafter\ifx\csname definition\endcsname\relax
\newenvironment{definition}{\begin{defi}}{\end{defi}}
\fi
\expandafter\ifx\csname lemma\endcsname\relax
\newenvironment{lemma}{\begin{lemm}}{\end{lemm}}
\fi
\expandafter\ifx\csname notation\endcsname\relax
\newenvironment{notation}{\begin{nota}}{\end{nota}}
\fi
\expandafter\ifx\csname problem\endcsname\relax
\newenvironment{problem}{\begin{pr}}{\end{pr}}
\fi
\expandafter\ifx\csname proposition\endcsname\relax
\newenvironment{proposition}{\begin{pro}}{\end{pro}}
\fi
\expandafter\ifx\csname remark\endcsname\relax
\newenvironment{remark}{\begin{rem}}{\end{rem}}
\fi
\expandafter\ifx\csname theorem\endcsname\relax
\newenvironment{theorem}{\begin{theo}}{\end{theo}}
\fi
\makeatother
\title[Van Vleck spectra of high-order Heun operators:\ finite-band]{Van Vleck spectra of high-order Heun operators:\ finite-band universality and exterior asymptotics}
\author{Boris Shapiro}
\date{Source: arXiv:2607.02700v1 (CC BY 4.0)}
\begin{document}
\maketitle

\noindent\emph{Source and licence.} Van Vleck spectra of high-order Heun operators:\ finite-band universality and exterior asymptotics. Version arXiv:2607.02700v1 (CC BY 4.0), distributed under the Creative Commons Attribution 4.0 International licence, as recorded in the arXiv licence metadata for this exact version and shown on the abstract page https://arxiv.org/abs/2607.02700v1. Full rights notice: licenses/arxiv-2607.02700-CC-BY-4.0.txt.\par\smallskip

\noindent\textit{Editorial scope.} Counted unit: the Proposition whose source label is \texttt{prop:k3-finiteband-pf} in the article (source0.tex lines 1112--1143, source label \texttt{prop:k3-finiteband-pf}), delivered together with the article's own complete original proof of it (source0.tex lines 1145--1276, one \texttt{proof} environment of 132 lines). No mathematics is written, repaired or re-derived: statement and proof are the article's own text line for line. Required context retained uncounted: none. Every definition, display and label of those lines is retained unchanged. Omitted from this package and not redistributed: 1--1111, 1144--1144, 1277--1675. Nothing inside a retained excerpt is omitted or altered: every retained body is delivered exactly as the article's lines are written, so no omission receipt of any kind accompanies this package. Citations: the retained excerpts carry no citation command at all, so no bibliography entry of the article is transcribed and no reference is redistributed. Reference adaptation: none was needed and none was made; every reference inside the delivered unit resolves inside it, so no unresolved reference or citation is produced. Printed slips kept literal: no printed word, symbol, hypothesis or number of the counted statement or of its proof was altered. Layout and class: the article's own class is \texttt{amsart} with packages including amsfonts, amsthm, latexsym, amsmath, amssymb, amscd, epsf, graphicx, float; the wrapper is the lane's pinned \texttt{amsart} editorial wrapper at 10pt on A4, which restates the theorem-like environments the retained text needs and adds the article's own macro definitions that the retained text needs (\texttt{\textbackslash CT}), with extra wrapper packages (bm, bbm, dsfont, xspace, cancel, mathtools). The delivered numbering is the unit's own. Assets: no retained span contains a figure environment, a graphics-inclusion command, a PDF-page inclusion, a table or a TikZ picture; no image file is copied into this package, the package declares no asset dependency and no image file is redistributed with it. Licence: version arXiv:2607.02700v1 of this article is distributed under the Creative Commons Attribution 4.0 International licence, as recorded in the arXiv licence metadata for this exact version and shown on the abstract page https://arxiv.org/abs/2607.02700v1. A full rights notice is delivered at licenses/arxiv-2607.02700-CC-BY-4.0.txt. Characters: every retained excerpt is pure ASCII, so the delivered engine, XeLaTeX with its default Latin Modern text font, reports no missing character.
\begin{proposition}[The cubic finite-band germ satisfies the cubic Picard--Fuchs equation]
\label{prop:k3-finiteband-pf}
Assume that the distinguished zero of the monic quartic leading coefficient has
been placed at the origin,
\[
        Q(t)=t^4+a_3t^3+a_2t^2+a_1t .
\]
Let \(C_Q(t)\) be the finite-band exterior germ
\[
        C_Q(t)=\int_0^1
        \frac{\Psi(t,\tau)\,d\tau}
        {3\Psi(t,\tau)^2-2A\Psi(t,\tau)-B},
\]
where
\[
        A=t+a_3\tau^3,\qquad
        B=(1-\tau^3)\tau^3a_2,\qquad
        D=(1-\tau^3)^2\tau^3a_1,
\]
and where \(\Psi=\Psi(t,\tau)\) is the branch satisfying
\(\Psi=t+O(1)\) at infinity of
\[
        \Psi^3=A\Psi^2+B\Psi+D .
\]
Then, as a germ at infinity, and hence by analytic continuation on the exterior
domain where \(C_Q\) is defined,
\begin{equation}
\label{eq:k3-finiteband-pf}
        81Q(t)C_Q'''(t)+162Q'(t)C_Q''(t)
        +90Q''(t)C_Q'(t)+10Q'''(t)C_Q(t)+30=0 .
\end{equation}
\end{proposition}

\begin{proof}
Put
\[
        b_\tau(w)=(1-\tau^3)w^{-1}
        -\tau^3(a_3+a_2w+a_1w^2).
\]
By the finite-band moment formula, the expansion of \(C_Q\) at infinity is
\[
        C_Q(t)=\frac1t+\sum_{m\ge1}\frac{M_m}{t^{m+1}},
        \qquad
        M_m=\int_0^1 \CT_w\, b_\tau(w)^m\,d\tau .
\]
We shall prove that these moments satisfy exactly the recurrence imposed by
\eqref{eq:k3-finiteband-pf}.

Expanding \(b_\tau(w)^m\), the constant term is obtained as follows.  Let
\(\ell\) be the number of factors \(-\tau^3a_2w\), and let \(r\) be the number
of factors \(-\tau^3a_1w^2\).  Then the number of factors
\((1-\tau^3)w^{-1}\) must be \(\ell+2r\), while the number of factors
\(-\tau^3a_3\) is \(m-2\ell-3r\).  Thus
\[
        M_m=\sum_{\substack{\ell,r\ge0\\2\ell+3r\le m}}
        c_m(\ell,r)\,
        a_3^{m-2\ell-3r}a_2^\ell a_1^r ,
\]
where
\[
        c_m(\ell,r)=
        \frac{(-1)^{m-\ell-2r}m!\,
        \Gamma(m-\ell-2r+\frac13)}
        {3\,\Gamma(m+\frac43)\,
        (m-2\ell-3r)!\,\ell!\,r!}.
\]
Indeed, after the substitution \(s=\tau^3\),
\[
        \int_0^1
        (1-\tau^3)^{\ell+2r}
        \tau^{3(m-\ell-2r)}\,d\tau
        =
        \frac13
        B\!\left(m-\ell-2r+\frac13,\ell+2r+1\right),
\]
which gives the displayed coefficient.

Now substitute
\[
        C_Q(t)=\sum_{m\ge0}M_m t^{-m-1},\qquad M_0=1,
\]
into the differential operator
\[
        \mathcal L
        =
        81Q(t)\frac{d^3}{dt^3}
        +162Q'(t)\frac{d^2}{dt^2}
        +90Q''(t)\frac{d}{dt}
        +10Q'''(t).
\]
The constant term of \(\mathcal L C_Q\) is \(-30\).  For \(m\ge1\), the
coefficient of \(t^{-m}\) in \(\mathcal L C_Q\) is \(-3R_m\), where
\[
\begin{split}
        R_m={}&
        (3m-5)(3m-2)(3m+1)M_m \\
        &+a_3(3m-5)(3m-2)^2M_{m-1} \\
        &+3a_2(m-1)(3m-5)(3m-4)M_{m-2} \\
        &+27a_1(m-2)^2(m-1)M_{m-3},
\end{split}
\]
with the convention \(M_j=0\) for \(j<0\).  Therefore it remains only to prove
\(R_m=0\) for every \(m\ge1\).

Because the \(M_m\)'s are polynomials in \(a_1,a_2,a_3\), it is enough to check
the coefficient of each monomial
\[
        a_3^{m-2\ell-3r}a_2^\ell a_1^r .
\]
The four possible contributions are
\[
        c_m(\ell,r),\qquad
        c_{m-1}(\ell,r),\qquad
        c_{m-2}(\ell-1,r),\qquad
        c_{m-3}(\ell,r-1),
\]
where inadmissible terms are understood as zero.  For admissible indices one has
\[
        \frac{c_{m-1}(\ell,r)}{c_m(\ell,r)}
        =
        -\frac{(m+\frac13)(m-2\ell-3r)}
        {m(m-\ell-2r-\frac23)},
\]
\[
        \frac{c_{m-2}(\ell-1,r)}{c_m(\ell,r)}
        =
        -\frac{\ell(m+\frac13)(m-\frac23)}
        {m(m-1)(m-\ell-2r-\frac23)},
\]
and
\[
        \frac{c_{m-3}(\ell,r-1)}{c_m(\ell,r)}
        =
        -\frac{r(m+\frac13)(m-\frac23)(m-\frac53)}
        {m(m-1)(m-2)(m-\ell-2r-\frac23)} .
\]
Substituting these three ratios into the coefficient of
\(a_3^{m-2\ell-3r}a_2^\ell a_1^r\) in \(R_m\), and bringing to a common
denominator, gives the identically zero numerator
\[
\begin{split}
& (3m-5)(3m-2)(3m+1) \\
&\quad
-(3m-5)(3m-2)^2
        \frac{(m+\frac13)(m-2\ell-3r)}
        {m(m-\ell-2r-\frac23)} \\
&\quad
-3\ell(m-1)(3m-5)(3m-4)
        \frac{(m+\frac13)(m-\frac23)}
        {m(m-1)(m-\ell-2r-\frac23)} \\
&\quad
-27r(m-2)^2(m-1)
        \frac{(m+\frac13)(m-\frac23)(m-\frac53)}
        {m(m-1)(m-2)(m-\ell-2r-\frac23)}
        =0 .
\end{split}
\]
Hence \(R_m=0\) for all \(m\ge1\).  Consequently
\[
        \mathcal L C_Q(t)+30=0
\]
as a Laurent series at infinity.  Since \(C_Q\) is analytic as an exterior germ,
the identity holds on its domain of analytic continuation.  This proves
\eqref{eq:k3-finiteband-pf}.
\end{proof}

\end{document}
